\section{Median of Two Sorted Arrays}
\label{sec:bs-median-two-arrays}

\problemstatement{We are given two sorted arrays $\textit{nums1}$ (size
$n_1$) and $\textit{nums2}$ (size $n_2$). We must find the median of the
combined sorted array of all $n_1 + n_2$ elements, \textbf{without}
actually merging the two arrays, in $O(\log(\min(n_1, n_2)))$ time.}

\begin{patternbox}
The keyword combination \emph{``median of two sorted arrays''} together
with a required complexity of $O(\log(\min(n_1,n_2)))$ -- far faster than
the $O(n_1+n_2)$ that an actual merge would cost -- is the signal for a
fundamentally different binary-search target than anything earlier in
this chapter: we will binary search not over values, and not over a
single array's indices, but over \textbf{how the combined sorted sequence
should be partitioned} between the two arrays. This ``binary search on a
partition point'' technique is one of the most intricate patterns on the
sheet, and is worth slowing down for.
\end{patternbox}

\subsection*{Understanding the problem carefully}

The median of a sorted sequence of length $T$ splits it into a left half
of size $\lceil T/2 \rceil$ and a right half of size $\lfloor T/2
\rfloor$, where every element in the left half is $\le$ every element in
the right half. Our goal is to find \emph{how many} elements of this left
half come from $\textit{nums1}$ (call this $\textit{cut1}$) and how many
come from $\textit{nums2}$ (call this $\textit{cut2} = \lceil T/2 \rceil -
\textit{cut1}$), \emph{without} ever materializing the merged array.

For a candidate $\textit{cut1} \in [0, n_1]$ (with $\textit{cut2}$
determined as above), define four boundary values:
\[
  l_1 = \textit{nums1}[\textit{cut1}-1] \ (\text{or } -\infty \text{ if }
  \textit{cut1}=0), \qquad
  r_1 = \textit{nums1}[\textit{cut1}] \ (\text{or } +\infty \text{ if }
  \textit{cut1}=n_1),
\]
\[
  l_2 = \textit{nums2}[\textit{cut2}-1] \ (\text{or } -\infty \text{ if }
  \textit{cut2}=0), \qquad
  r_2 = \textit{nums2}[\textit{cut2}] \ (\text{or } +\infty \text{ if }
  \textit{cut2}=n_2).
\]
A partition is \textbf{valid} exactly when $l_1 \le r_2$ \emph{and}
$l_2 \le r_1$: every element on the left side ($l_1$ and $l_2$) is at most
every element on the right side ($r_1$ and $r_2$), which is precisely the
condition needed for the left half (across both arrays combined) to
contain only values $\le$ every value in the right half.

\begin{ideabox}[Candidate Space and Predicate]
The candidate space is $\textit{cut1} \in [0, n_1]$ (binary searching on
the \emph{smaller} array keeps the range, and hence the iteration count,
as small as possible). The predicate is
$P(\textit{cut1}) = (l_1 \le r_2 \text{ and } l_2 \le r_1)$. Because
increasing $\textit{cut1}$ monotonically increases $l_1$ and $r_1$ while
monotonically decreasing $l_2$ and $r_2$ (as $\textit{cut2}$ shrinks
correspondingly), there is exactly one crossover region where the
partition is valid, and binary search finds it directly.
\end{ideabox}

\subsection*{Why exactly one valid partition exists, and how to search for it}

If $l_1 > r_2$, the left part of $\textit{nums1}$ is reaching too far
(taking values larger than what the right part of $\textit{nums2}$ can
absorb), so $\textit{cut1}$ must decrease: search left,
$\textit{hi} \gets \textit{cut1}-1$. If $l_2 > r_1$ (equivalently,
$\textit{nums1}$'s left part is not reaching far enough), $\textit{cut1}$
must increase: search right, $\textit{lo} \gets \textit{cut1}+1$. Because
$l_1, r_1$ increase weakly with $\textit{cut1}$ and $l_2, r_2$ decrease
weakly with $\textit{cut1}$ (equivalently increase weakly as
$\textit{cut1}$ decreases), exactly one of these two failure conditions
can hold at a time for an invalid partition, and the search converges to
the unique boundary where neither failure condition holds.

Once a valid partition is found, the median is computed from
$\max(l_1, l_2)$ and $\min(r_1, r_2)$: if $T$ is odd, the median is
$\max(l_1, l_2)$ (the largest element of the left half, which has one more
element than the right half); if $T$ is even, the median is the average of
$\max(l_1, l_2)$ and $\min(r_1, r_2)$.

\subsection*{Algorithm}

\begin{enumerate}
  \item Ensure $\textit{nums1}$ is the smaller array (swap if necessary,
        so the binary search range is $[0, n_1]$ with $n_1 \le n_2$).
  \item $T \gets n_1+n_2$, $\textit{half} \gets \lceil T/2 \rceil$.
  \item $\textit{lo} \gets 0$, $\textit{hi} \gets n_1$.
  \item While $\textit{lo} \le \textit{hi}$:
  \begin{enumerate}
    \item $\textit{cut1} \gets \textit{lo}+(\textit{hi}-\textit{lo})/2$,
          $\textit{cut2} \gets \textit{half}-\textit{cut1}$.
    \item Compute $l_1, r_1, l_2, r_2$ as above.
    \item If $l_1 \le r_2$ and $l_2 \le r_1$: valid partition found;
          compute and return the median as described above.
    \item Else if $l_1 > r_2$: $\textit{hi} \gets \textit{cut1}-1$.
    \item Else: $\textit{lo} \gets \textit{cut1}+1$.
  \end{enumerate}
\end{enumerate}

\begin{algorithm}[H]
\caption{Median of Two Sorted Arrays}
\begin{algorithmic}[1]
\Procedure{FindMedian}{$\textit{nums1}, \textit{nums2}$}
  \If{$n_1 > n_2$} \textbf{swap} $\textit{nums1}, \textit{nums2}$ \EndIf
  \State $T \gets n_1 + n_2$, $\textit{half} \gets \lceil T/2 \rceil$
  \State $\textit{lo} \gets 0$, $\textit{hi} \gets n_1$
  \While{$\textit{lo} \le \textit{hi}$}
    \State $\textit{cut1} \gets \textit{lo}+(\textit{hi}-\textit{lo})/2$
    \State $\textit{cut2} \gets \textit{half} - \textit{cut1}$
    \State $l_1 \gets (\textit{cut1}=0)\ ? -\infty : \textit{nums1}[\textit{cut1}-1]$
    \State $r_1 \gets (\textit{cut1}=n_1)\ ? +\infty : \textit{nums1}[\textit{cut1}]$
    \State $l_2 \gets (\textit{cut2}=0)\ ? -\infty : \textit{nums2}[\textit{cut2}-1]$
    \State $r_2 \gets (\textit{cut2}=n_2)\ ? +\infty : \textit{nums2}[\textit{cut2}]$
    \If{$l_1 \le r_2$ \textbf{and} $l_2 \le r_1$}
      \If{$T$ is odd} \Return $\max(l_1, l_2)$
      \Else \Return $(\max(l_1,l_2)+\min(r_1,r_2))/2$ \EndIf
    \ElsIf{$l_1 > r_2$}
      \State $\textit{hi} \gets \textit{cut1} - 1$
    \Else
      \State $\textit{lo} \gets \textit{cut1} + 1$
    \EndIf
  \EndWhile
\EndProcedure
\end{algorithmic}
\end{algorithm}

\subsection*{Dry run}

\dryrun{Take $\textit{nums1} = [1,3]$, $\textit{nums2} = [2]$.}

$n_1=2 \le n_2=1$? No -- swap so $\textit{nums1}=[2]$ ($n_1=1$),
$\textit{nums2}=[1,3]$ ($n_2=2$). $T=3$, $\textit{half}=\lceil 3/2
\rceil=2$.

\begin{itemize}
  \item $\textit{lo}=0,\textit{hi}=1$: $\textit{cut1}=0$,
        $\textit{cut2}=2-0=2$. $l_1=-\infty$ ($\textit{cut1}=0$),
        $r_1=\textit{nums1}[0]=2$. $l_2=\textit{nums2}[1]=3$
        ($\textit{cut2}=2=n_2$, so $r_2=+\infty$). Check: $l_1=-\infty \le
        r_2=+\infty$? Yes. $l_2=3 \le r_1=2$? No. So $l_2 > r_1$,
        $\textit{lo} \gets 1$.
  \item $\textit{lo}=1,\textit{hi}=1$: $\textit{cut1}=1$,
        $\textit{cut2}=2-1=1$. $l_1=\textit{nums1}[0]=2$
        ($\textit{cut1}=1=n_1$, so $r_1=+\infty$). $l_2=\textit{nums2}[0]=1$,
        $r_2=\textit{nums2}[1]=3$. Check: $l_1=2 \le r_2=3$? Yes.
        $l_2=1 \le r_1=+\infty$? Yes. Valid partition.
\end{itemize}

$T=3$ is odd, so the median is $\max(l_1,l_2) = \max(2,1) = 2$. Checking
by hand: the merged array is $[1,2,3]$, whose median is indeed $2$.

\subsection*{C++ implementation}

\begin{lstlisting}
#include <bits/stdc++.h>
using namespace std;

double findMedianSortedArrays(vector<int>& nums1, vector<int>& nums2) {
    if (nums1.size() > nums2.size()) swap(nums1, nums2);

    int n1 = nums1.size(), n2 = nums2.size();
    int total = n1 + n2;
    int half = (total + 1) / 2;

    int lo = 0, hi = n1;

    while (lo <= hi) {
        int cut1 = lo + (hi - lo) / 2;
        int cut2 = half - cut1;

        int l1 = (cut1 == 0) ? INT_MIN : nums1[cut1 - 1];
        int r1 = (cut1 == n1) ? INT_MAX : nums1[cut1];
        int l2 = (cut2 == 0) ? INT_MIN : nums2[cut2 - 1];
        int r2 = (cut2 == n2) ? INT_MAX : nums2[cut2];

        if (l1 <= r2 && l2 <= r1) {
            if (total % 2 == 1) {
                return max(l1, l2);
            } else {
                return (max(l1, l2) + min(r1, r2)) / 2.0;
            }
        } else if (l1 > r2) {
            hi = cut1 - 1;
        } else {
            lo = cut1 + 1;
        }
    }
    return -1.0; // unreachable for valid sorted input
}

int main() {
    vector<int> nums1 = {1, 3};
    vector<int> nums2 = {2};
    cout << "Median: "
         << findMedianSortedArrays(nums1, nums2) << endl;
    return 0;
}
\end{lstlisting}

\begin{complexitybox}
\textbf{Time Complexity.} The binary search runs over $\textit{cut1} \in
[0, n_1]$, where $n_1 = \min(n_1, n_2)$ (after the swap), giving
\[
  O(\log(\min(n_1, n_2))).
\]
\textbf{Space Complexity.} $O(1)$ auxiliary space.
\end{complexitybox}

\begin{takeawaybox}
Binary search does not need to search over array \emph{values} or even a
single array's \emph{indices} -- here, it searches over a \textbf{split
point} that simultaneously partitions two arrays, using a validity
condition ($l_1 \le r_2$ and $l_2 \le r_1$) as the monotonic predicate.
This ``binary search on a partition'' idea, once recognized, also unlocks
Section~\ref{sec:bs-kth-element}, the general $k$-th element version of
this same problem.
\end{takeawaybox}
